\documentclass[12pt]{report}
\usepackage[english]{babel}
\usepackage[utf8x]{inputenc}
\usepackage{amsmath}
\usepackage{graphicx}
\usepackage[colorinlistoftodos]{todonotes}
\usepackage[left=2cm,right=2cm,top=2.5cm,bottom=2cm]{geometry}
\usepackage[]{algorithm2e}
\usepackage{amsmath}
\usepackage{amssymb}

\begin{document}

%----------------------------------------------------------------------------------------
%	MAIN PAGE
%----------------------------------------------------------------------------------------
\begin{titlepage}
\vspace*{\fill}
\newcommand{\HRule}{\rule{\linewidth}{0.5mm}}

\center

\textsc{\LARGE Aix-Marseille Université}\\[1.5cm]
\textsc{\Large Compléxité}\\[0.5cm]

\HRule \\[0.4cm]
{ \huge \bfseries Deux réductions vers SAT}\\[0.4cm] % Title of your document
\HRule \\[1.5cm]

\begin{minipage}{0.4\textwidth}
\begin{flushleft} \large
\emph{Autheur:}\\
Michaël \textsc{Gileta} \\
Yohan \textsc{Roux}
\end{flushleft}
\end{minipage}
~
\begin{minipage}{0.4\textwidth}
\begin{flushright} \large
\emph{Référent:} \\
Kévin \textsc{Perrot}
\end{flushright}
\end{minipage}\\[2cm]

{\large Vendredi 18 Octobre}\\[2cm]

\includegraphics[scale=0.5]{./images/logo.png}\\[1cm]

\vspace*{\fill}
\end{titlepage}

%----------------------------------------------------------------------------------------

\section*{Réduction de KERNEL à SAT}
\subsection*{1. Réduction}

Soit s un sommet, avec pour tous s $\in$ V, si $s=1$ $\Leftrightarrow$ s $\in$ $V\prime$, sinon $ s=0 \Leftrightarrow$ s $\in$ V. \\
Soit $(s\prime, s\prime\prime)$ un arc orienté de $s\prime$ vers $s\prime\prime$\\
\\
Pour un graph G=(V,E) avec : \\
V : l'ensemble des sommets\\
E : l'ensemble des arc orienté\\
\\
On construit notre formule à partir de E puis de V\\
\\
(a) $\underset{C_{ij} \in E}{\land} (\lnot i \lor \lnot j)$ \\
(b) $\underset{j \in V}{\land} (\lnot j \underset{C_{ij} \in E}{\lor} (\lnot i))$

\subsection*{2. Minisat}
\section*{Réduction de SUDOKU à SAT}
\subsection*{1. Modélisation de contrainte}
Soit E = \{1..9\} , $E\prime$=\{0,1,2\} , $E\prime\prime$=\{1,2,3\}\\
Soit $C_{i,j,k} \in E$\\
La case (i,j) contient la valeur k
\subsubsection*{Au moins une valeur par case}
$\underset{i,j \in E}{\land} \underset{k \in E}{\lor} C_{ijk}$
\subsubsection*{Au moins une fois chaque chiffre sur chaque ligne}
$\underset{i,k \in E}{\land} \underset{j \in E}{\lor} C_{ijk}$
\subsubsection*{Au moins une fois chaque chiffre sur chaque colonne}
$\underset{j,k \in E}{\land} \underset{i \in E}{\lor} C_{ijk}$
\subsubsection*{Au moins une fois chaque chiffre dans chaque region}
$\underset{k \in E}{\underset{i,j \in E\prime}{\land}} \underset{i\prime,j\prime \in E\prime\prime,}{\land} (C_{[(3*i)+i\prime] [(3*j)+j\prime]k})$
\subsubsection*{Au plus une valeur par case}
$\underset{i,j \in E}{\land} \underset{\underset{k \ne k\prime}{k,k\prime \in E,} }{\lor} (\lnot C_{ijk} \lor \lnot C_{ijk\prime})$
\subsubsection*{Au plus une fois chaque chiffre sur chaque ligne}
$\underset{i,k \in E}{\land} \underset{\underset{j \ne j\prime}{j,j\prime \in E,}}{\lor} (\lnot C_{ijk} \lor \lnot C_{ij\prime k})$
\subsubsection*{Au plus une fois chaque chiffre sur chaque colonne}
$\underset{j,k \in E}{\land} \underset{\underset{i \ne i\prime}{i,i\prime \in i,}}{\lor} (\lnot C_{ijk} \lor \lnot C_{i\prime jk})$
\subsubsection*{Au plus une fois chaque chiffre dans chaque region}
$\underset{i,j,k \in E}{\land} \underset{\underset{j \neq j\prime}{\underset{i \neq i\prime,}{\underset{j\prime \in E\prime\prime,}{i\prime \in E\prime\prime,}}}}{\land} (\lnot C_{ijk} \lor \lnot C_{[(3*i)+i\prime] [(3*j)+j\prime]k})$

\end{document}